\clearpage
\section{RX-side Gain Flatness}
\label{sec:rx-flatness}
\raggedbottom% no stretched gaps between text and figures (as in the NF section)

\textbf{What is gain flatness?} Exactly as on the TX side (Section~\ref{sec:tx-flatness}): how much the conversion
gain varies \emph{within} a 500\,MHz window, rather than its average (Sections~\ref{sec:rx-gain-lowside} and
\ref{sec:rx-gain-highside}). A received wideband signal that passes through an uneven channel comes out of the unit
with a tilted or rippled spectrum.

\textbf{How it's measured.} The same way as on the TX side, with the roles of the ports swapped: at every (FR1, FR3)
point, the FieldFox's tracking generator sweeps the FR3 window into the ``RX-side FR3 RF In'' port, while its
receiver follows the 500\,MHz window centred on FR1 at the ``RX-side FR1 IF Out'' port, in 101 steps 5\,MHz apart.
The input level is −50\,dBm, the same as in the RX conversion-gain sections, low enough to keep the highest-gain
points clear of compression, which would flatten the measured ripple. The signal level arriving at the
unit was again measured directly, by connecting the FieldFox's source to its own receiver through the same input
cable, and the output cable's loss is removed from its own calibration. Each window is measured with the AM3153
filter bypassed and with the ``optimal'' filter, never changing the filter during a sweep, over the same grid as on
the TX side, for both low-side and high-side LO. The number reported is the RMS flatness defined in
Section~\ref{sec:tx-flatness}: 0\,dB is perfectly flat, and in this data the peak-to-peak spread in a window is
typically about \flatRxPpRatio{} times the RMS flatness (between \flatRxPpRatioLow{} and \flatRxPpRatioHigh{} times
at 80\% of the points).

\textbf{Measurement conditions.} \lmxDriveRule{} (set from the LO frequency). Because the LO is shared, the TX side was set up at every point exactly as in normal use (its stored LO-leakage calibration applied and its optimal filter engaged), with both TX-side ports terminated in 50\,\textohm. The unit was measured hot, as on the TX side: no fan, LMX2820 die at \condRxFlatMinC--\condRxFlatMaxC\textdegree C.

\textbf{Results.} Within the design spec (FR3~$\le$~24\,GHz, and FR1 centre~$\le$~3.5\,GHz):
\par\medskip\noindent\begin{minipage}{\linewidth}% caption and table stay on one page
\begin{center}
\captionof{table}{RX-side RMS gain flatness within the design spec, by LO side and filter setting.}\label{tab:rxflat-results}
\vspace{0.4em}
\begin{tabular}{lcccc}
\toprule
 & Median & 90\% of points at or below & Best & Worst \\
\midrule
Low-side LO, filter bypassed & \flatRxLowBypassMedian\,dB & \flatRxLowBypassPninety\,dB & \flatRxLowBypassMin\,dB & \flatRxLowBypassMax\,dB \\
Low-side LO, optimal filter & \flatRxLowOptimalMedian\,dB & \flatRxLowOptimalPninety\,dB & \flatRxLowOptimalMin\,dB & \flatRxLowOptimalMax\,dB \\
High-side LO, filter bypassed & \flatRxHighBypassMedian\,dB & \flatRxHighBypassPninety\,dB & \flatRxHighBypassMin\,dB & \flatRxHighBypassMax\,dB \\
High-side LO, optimal filter & \flatRxHighOptimalMedian\,dB & \flatRxHighOptimalPninety\,dB & \flatRxHighOptimalMin\,dB & \flatRxHighOptimalMax\,dB \\
\bottomrule
\end{tabular}
\end{center}
\end{minipage}\par

\textbf{Where it's flattest, and how to use that.} The median per FR1 centre (all FR3 up to 24\,GHz):
\par\medskip\noindent\begin{minipage}{\linewidth}% caption and table stay on one page
\begin{center}
\captionof{table}{RX-side median RMS gain flatness for each FR1 centre (all FR3 up to 24\,GHz).}\label{tab:rxflat-by-fr1}
\vspace{0.4em}
\footnotesize
% Generated by FR3_Characterization/plot_datasheet_flatness.py -- do not edit by hand.
\begin{tabular}{lcccc}
\toprule
FR1 centre & Low-side LO, bypassed & Low-side LO, optimal & High-side LO, bypassed & High-side LO, optimal \\
\midrule
1.0\,GHz & 0.91\,dB & 0.74\,dB & 0.85\,dB & 0.77\,dB \\
1.5\,GHz & 1.75\,dB & 1.30\,dB & 1.64\,dB & 1.29\,dB \\
2.0\,GHz & 1.11\,dB & 1.00\,dB & 1.03\,dB & 0.86\,dB \\
2.5\,GHz & 0.79\,dB & 0.75\,dB & 0.71\,dB & 0.76\,dB \\
3.0\,GHz & 0.70\,dB & 0.81\,dB & 0.64\,dB & 0.71\,dB \\
3.5\,GHz & 1.14\,dB & 1.20\,dB & 1.04\,dB & 1.64\,dB \\
4.0\,GHz\textsuperscript{*} & 1.06\,dB & 0.92\,dB & 0.90\,dB & 1.31\,dB \\
\bottomrule
\multicolumn{5}{l}{\footnotesize\textsuperscript{*}Its 500\,MHz window (3.75--4.25\,GHz) runs 250\,MHz past the IF design spec.}
\end{tabular}

\end{center}
\end{minipage}\par
The RX side is flatter than the TX side (Section~\ref{sec:tx-flatness}): its median RMS flatness is about 1\,dB
or less in every configuration. As on the TX side, it depends mostly on where the signal sits in FR1, and much less on
FR3 or on the LO side. Windows centred at FR1 = 2.5--3.0\,GHz are the flattest, closely followed by 1.0\,GHz, while
windows centred at 1.5\,GHz are the least flat, followed by 3.5\,GHz. So when flatness matters and you are free to
choose the IF, centre the signal near 2.5--3.0\,GHz in FR1 and avoid 1.5\,GHz. Engaging the optimal filter changes
the flatness only a little: its median is slightly lower than with the filter bypassed, for both LO sides.

\textbf{What's plotted.} As on the TX side: for each LO side and filter setting, a heatmap of the RMS flatness over
the full measured (FR1, FR3) grid, on the shared 0--4\,dB colour scale (darker = flatter), with points above 24\,GHz
and the FR1 = 4.0\,GHz column boxed in red as out of design spec; and a histogram of the points within the design spec.

\figpair{images/flat_rx_low_bypass_map.pdf}{RMS gain flatness, low-side LO: full measured (FR1, FR3) grid, filter
bypassed. Zoom in to read each cell's exact values, printed as (FR1, FR3, RMS flatness), with FR1 and FR3 in GHz and
the RMS flatness in dB.}{fig:rxflat-1}%
        {images/flat_rx_low_bypass_hist.pdf}{LO: Low-side. Filter: Bypassed. Distribution of
Figure~\ref{fig:rxflat-1}, points within the design spec.}{fig:rxflat-2}
\figpair{images/flat_rx_low_optimal_map.pdf}{RMS gain flatness, low-side LO: full measured (FR1, FR3) grid,
``optimal'' filter engaged at every point.}{fig:rxflat-3}%
        {images/flat_rx_low_optimal_hist.pdf}{LO: Low-side. Filter: Optimal. Distribution of
Figure~\ref{fig:rxflat-3}, points within the design spec.}{fig:rxflat-4}
\clearpage
\figpair{images/flat_rx_high_bypass_map.pdf}{RMS gain flatness, high-side LO: full measured (FR1, FR3) grid, filter
bypassed.}{fig:rxflat-5}%
        {images/flat_rx_high_bypass_hist.pdf}{LO: High-side. Filter: Bypassed. Distribution of
Figure~\ref{fig:rxflat-5}, points within the design spec.}{fig:rxflat-6}
\figpair{images/flat_rx_high_optimal_map.pdf}{RMS gain flatness, high-side LO: full measured (FR1, FR3) grid,
``optimal'' filter engaged at every point.}{fig:rxflat-7}%
        {images/flat_rx_high_optimal_hist.pdf}{LO: High-side. Filter: Optimal. Distribution of
Figure~\ref{fig:rxflat-7}, points within the design spec.}{fig:rxflat-8}
\flushbottom
